\begin{tabular}{lrrr}
\toprule
\textbf{Punto} & \textbf{Cs/Co con 12 ns} & \textbf{Cs/Co con 75 ns} & \textbf{Cambio [\%]} \\
\midrule
A @ 1600 V & 1.0750 & 1.0884 & $+1.24$ \\
A @ 1650 V & 1.1033 & 1.1092 & $+0.53$ \\
A @ 1700 V & 1.1238 & 1.1337 & $+0.88$ \\
B @ 1600 V & 1.0682 & 1.0796 & $+1.07$ \\
B @ 1650 V & 1.0663 & 1.0748 & $+0.80$ \\
B @ 1700 V & 1.0765 & 1.0756 & $-0.08$ \\
C @ 1600 V & 0.8693 & 0.8762 & $+0.79$ \\
C @ 1650 V & 0.9015 & 0.9100 & $+0.93$ \\
C @ 1700 V & 0.8908 & 0.9013 & $+1.18$ \\
\bottomrule
\end{tabular}
\par\smallskip
\begin{minipage}{0.95\linewidth}\footnotesize Donde el registro alcanza para las dos compuertas, el cociente entre fuentes se mueve menos del 1\%: la eleccion de la compuerta no explica el desacuerdo del test A. En el detector D el control no se puede correr, porque con 75 ns entra el 2\% de los eventos.\end{minipage}
